%---------------------------------------------------------------
% Resume in LaTeX -- Ryan Strawbridge
% Template originally by Vishhal Shashi Kumar.
%
% Build:  tectonic ryan-strawbridge-resume.tex
% Output: ryan-strawbridge-resume.pdf
%
% If it spills onto page 2, shave \entrygap and \sectiongap first
% (1pt at a time). If it looks too airy, do the opposite.
%---------------------------------------------------------------

\documentclass[letterpaper,10pt]{article}

\usepackage{latexsym}
\usepackage[empty]{fullpage}
\usepackage{titlesec}
\usepackage{marvosym}
\usepackage[usenames,dvipsnames]{color}
\usepackage{verbatim}
\usepackage{enumitem}
\usepackage[hidelinks]{hyperref}
\usepackage{fancyhdr}
\usepackage[english]{babel}
\usepackage{tabularx}
\usepackage{mathpazo}
\usepackage[T1]{fontenc}

% ATS-friendly glyph mapping. These are pdfTeX primitives; XeTeX (what
% tectonic runs) doesn't have them, so guard it and the file compiles under
% both pdflatex (Overleaf) and tectonic.
\ifdefined\pdfgentounicode
  \input{glyphtounicode}
\fi

% slightly open leading -- Palatino/Charter have a large x-height
\linespread{1.05}

\pagestyle{fancy}
\fancyhf{}
\fancyfoot{}
\renewcommand{\headrulewidth}{0pt}
\renewcommand{\footrulewidth}{0pt}

%---------------------------------------------------------------
% Margins
%---------------------------------------------------------------

\addtolength{\oddsidemargin}{-0.5in}
\addtolength{\evensidemargin}{-0.5in}
\addtolength{\textwidth}{1in}
\addtolength{\topmargin}{-0.65in}
\addtolength{\textheight}{1.45in}

\urlstyle{same}
\raggedbottom
\raggedright
\setlength{\tabcolsep}{0in}

%===============================================================
% SPACING CONTROLS -- the only numbers you should need to touch
%===============================================================

\newlength{\sectiongap}   \setlength{\sectiongap}{8pt}   % above each section title
\newlength{\afterrulegap} \setlength{\afterrulegap}{3pt} % below the section rule
\newlength{\entrygap}     \setlength{\entrygap}{4pt}     % between jobs / entries
\newlength{\bulletgap}    \setlength{\bulletgap}{1pt}    % between bullets in one entry
\newlength{\headgap}      \setlength{\headgap}{2pt}      % heading block to its bullets
\newlength{\rolegap}      \setlength{\rolegap}{13pt}     % pull-up under a second-role line

%---------------------------------------------------------------
% Section formatting
%---------------------------------------------------------------

\titleformat{\section}
  {\scshape\raggedright\large}{}{0em}{}
  [\color{black}\titlerule]
\titlespacing*{\section}{0pt}{\sectiongap}{\afterrulegap}

% ATS friendly (pdfTeX only -- see the guard above)
\ifdefined\pdfgentounicode
  \pdfgentounicode=1
\fi

%---------------------------------------------------------------
% Custom commands
%---------------------------------------------------------------

\newcommand{\resumeItem}[1]{\item\small{#1}}

\newcommand{\resumeSubheading}[4]{
  \vspace{\entrygap}\item
    \begin{tabular*}{\textwidth}[t]{l@{\extracolsep{\fill}}r}
      \textbf{#1} & #2 \\
      \textit{\small #3} & \textit{\small #4}
    \end{tabular*}\vspace{-5pt}
}

% Second role at the same employer -- no company line repeated.
% This tabular has ONE row where \resumeSubheading has two. With [t] the box
% hangs below the baseline by its own height, so the -5pt that suits the
% two-row version leaves a visible gap here. \rolegap is the pull-up; raise it
% if the bullets sit too far below the role line, lower it if they collide.
\newcommand{\resumeSubSubheading}[2]{
  \vspace{2pt}\item
    \begin{tabular*}{\textwidth}[t]{l@{\extracolsep{\fill}}r}
      \textit{\small #1} & \textit{\small #2}
    \end{tabular*}\vspace{-\rolegap}
}

% Two-line education / degree entry
\newcommand{\resumeEducation}[5]{
  \vspace{\entrygap}\item
    \begin{tabular*}{\textwidth}[t]{l@{\extracolsep{\fill}}r}
      \textbf{#1} \textit{\small -- #2} & #3 \\
      \textit{\small #4} & \textit{\small #5}
    \end{tabular*}\vspace{-4pt}
}

\newcommand{\resumeProjectHeading}[2]{
  \vspace{\entrygap}\item
    \begin{tabular*}{\textwidth}[t]{l@{\extracolsep{\fill}}r}
      \small #1 & \small #2
    \end{tabular*}\vspace{-8pt}
}

\renewcommand\labelitemii{$\vcenter{\hbox{\tiny$\bullet$}}$}

\newcommand{\resumeSubHeadingListStart}{
  \begin{itemize}[
    leftmargin=0in,
    label={},
    topsep=-\entrygap,
    itemsep=0pt,
    parsep=0pt,
    partopsep=0pt
  ]
}

\newcommand{\resumeSubHeadingListEnd}{\end{itemize}}

\newcommand{\resumeItemListStart}{
  \begin{itemize}[
    leftmargin=0.20in,
    topsep=\headgap,
    itemsep=\bulletgap,
    parsep=0pt,
    partopsep=0pt
  ]
}

\newcommand{\resumeItemListEnd}{\end{itemize}\vspace{-5pt}}

%===============================================================
% RESUME
%===============================================================

\begin{document}

%----------HEADING----------

\begin{center}
    \textbf{\LARGE Ryan Strawbridge} \\ \vspace{3pt}
    \small
    508.594.9661 $|$
    \href{mailto:strawbridge.r@northeastern.edu}
    {\underline{strawbridge.r@northeastern.edu}} $|$
    \href{https://www.linkedin.com/in/ryan-strawbridge}
    {\underline{LinkedIn}} $|$
    \href{https://ryanstrawbridge-2.github.io}
    {\underline{Portfolio}} $|$
    Boston, MA
\end{center}

\vspace{-10pt}

%-----------EDUCATION-----------

\section{Education}

\resumeSubHeadingListStart

  \resumeEducation
    {Northeastern University}{Boston, MA}{May 2027}
    {Bachelor of Science in Mechanical Engineering (BSME)}{GPA: 3.58}

\resumeSubHeadingListEnd

\vspace{\entrygap}

{\small
\textbf{Coursework:}
Heat Transfer | Fluid Mechanics | Solid Mechanics |
Thermodynamics | Material Science | Dynamics \\[3pt]
\textbf{Activities:}
AeroNU Propulsion Member $|$ Running Club President $|$
Rhythm Guitarist in Alt. Rock Band \\[3pt]
\textbf{Achievements:}
DoD Security Clearance $|$ CSWA Certified $|$ Lean Six Sigma Certification $|$
800th in the Boston Marathon
}

%-----------PROFESSIONAL EXPERIENCE-----------

\section{Professional Experience}

\resumeSubHeadingListStart

% CFS -- TERM 2
\resumeSubheading
  {Commonwealth Fusion Systems}
  {Devens, MA}
  {Manufacturing / Equipment Engineering Co-op (Term 2)}
  {Jan 2026 -- Jun 2026}

\resumeItemListStart

  \resumeItem{
  Designed and validated a 2,100 lb hand-portable lifting frame to
  \textbf{BTH-1} and \textbf{ASME B30.20}, earning sign-off from structural
  engineering and safety; increased table space for fixture tooling during
  magnet stacking and \textbf{saved 30--45 minutes per lift}}

  \resumeItem{
  Developed a modular CNC seam-cutting fixture accommodating 11 cable
  configurations, enabling qualification testing ahead of a new production
  line; \textbf{applied GD\&T} with position and symmetry controls on dowel-pin datums
  for repeatable alignment, \textbf{reducing future operation time by over 80\%}}

  \resumeItem{
  Deployed a low-cost adjustable fixture using welded tabs and off-the-shelf
  linear rails and carriages to maintain cable alignment during insulation
  wrapping, improving process robustness and \textbf{reducing labor by 0.5--3 hours
  per magnet layer}}

\resumeItemListEnd

% CFS -- TERM 1
\resumeSubSubheading
  {Manufacturing Engineering Co-op (Term 1)}
  {Jan 2025 -- Jun 2025}

\resumeItemListStart

  \resumeItem{
  Re-designed a hydraulic piston-driven bending machine in Siemens NX using
  bolt safety analysis and first-principles calculations, cutting operation
  time by 75\% and reducing injury risk; \textbf{still in use over a year later,
  saving hundreds of labor hours}}

  \resumeItem{
  Developed a predictive equation relating bend offset, bend radius, and arc
  length through iterative spring-back experiments and data interpolation,
  \textbf{eliminating trial-and-error die design} and enabling future cable-bending
  dies to be sized analytically on the first iteration}

  \resumeItem{
  \textbf{Owned a work cell} and wrote detailed work instructions with technicians
  covering cable assembly, joint placement, orbital cutting, and welding,
  replacing tribal knowledge with a repeatable reference used on the
  production floor}

\resumeItemListEnd

% RAYTHEON
\resumeSubheading
  {Raytheon Missiles \& Defense}
  {Andover, MA}
  {Operations Engineering Co-op}
  {Jan 2024 -- Sep 2024}

\resumeItemListStart

  \resumeItem{
  Prototyped a composite fixture using SolidWorks and PCB Gerber files to
  shield heat-sensitive filters from 93\,\textdegree C while enabling reflow of
  adjacent components, \textbf{increasing first-pass yield} and removing filter damage
  as a scrap driver on the line}

  \resumeItem{
  Ran a nine-configuration solder stencil study varying placement and volume
  to provide off-gas escape channels, reducing voiding and \textbf{removing a rework
  step} from the process}

  \resumeItem{
  \textbf{Co-owned a work cell} and developed thermal reflow profiles by refining oven
  parameters through thermocouple data, operator input, and X-ray inspection
  to IPC-J-STD-001 Class 3, \textbf{reducing rework time across 15 programs}}

\resumeItemListEnd

\resumeSubHeadingListEnd

%-----------PROJECTS-----------

\section{Projects}

\resumeSubHeadingListStart

% CAPSTONE
\resumeSubheading
  {Underwater Acoustic Release Mechanism -- Capstone Project}
  {Boston, MA}
  {Mechanical Design Lead -- Release System}
  {Jun 2026 -- Present}

\resumeItemListStart

  \resumeItem{
  \textbf{Own the mechanical design} of a low-cost underwater acoustic release
  mechanism for underwater robot recovery; designing a rotary lever arm
  mechanism and using first-principles force analysis to evaluate and rule out
  initial design concepts}

  \resumeItem{
  Performed \textbf{structural analysis} of the release shaft under combined axial
  compression and torsion, and selected O-ring seal geometry using Parker's
  O-Ring Handbook to keep the rotating, pressure-sealed shaft interface
  leak-free at 300 m depth}

  \resumeItem{
  Validating hand-calculated stress predictions on housing and latch
  components in \textbf{Ansys FEA} under combined pressure and mechanical loading}

\resumeItemListEnd

% AERONU
\resumeSubheading
  {AerospaceNU -- Liquid Rocketry Club}
  {Boston, MA}
  {Liquid Rocket Engine Member}
  {Sep 2023 -- Present}

\resumeItemListStart

  \resumeItem{
  Designing and validating an actively controlled nitrogen pressurant
  subsystem, integrating a servo-actuated ball valve and torque coupling with
  closed-loop pressure regulation based on downstream pressure feedback}

  \resumeItem{
  Supported setup and execution of \textbf{three successful static-fire tests},
  integrating propulsion, electrical, and instrumentation systems while
  configuring test hardware and data acquisition equipment}

\resumeItemListEnd

\resumeSubHeadingListEnd

%-----------SKILLS-----------

\section{Technical Skills}

\vspace{6.5pt}
{\small
\textbf{Software:}
Siemens NX | SolidWorks | Ansys Mechanical | Arduino | Teamcenter | ManufacturO | Odoo | Confluence | Jira \\[3pt]
\textbf{Design:}
Structural Analysis, DFM, GD\&T, FEA (Ansys -- developing) \\[3pt]
\textbf{Fabrication:}
Manual G-Code | CNC Milling | Power Tools | TIG Welding | 3D Printing | Soldering | Water Jetting | Laser Cutting
}

%---------------------------------------------------------------

\end{document}
